% GENERATED FILE. Edit canonical lessons, then run npm run generate:books.
% canonical-source-hash: fnv1a64:f69b3669064b6d0f
% canonical-lessons: PT-C528-aspirar, PT-C528-aspirador, PT-C528-varrer, PT-C528-esfregar, PT-C528-esfregona

\chapter{Dia de limpezas}
\label{ch:pt-dia-de-limpezas}
\hlchaptermodality{528}

\begin{chapteropening}
\textbf{By the end of this chapter:} \emph{I can talk about cleaning the house.}

Complete the last lesson of chapter 528: I can talk about cleaning the house.
\end{chapteropening}

\section[aspirar]{aspirar — to vacuum, to hoover; to breathe in}

\label{lesson:PT-C528-aspirar}



\begin{quote}
Before the new one: say the Portuguese for documentary, then the Portuguese for cast.
\end{quote}

\subsection*{You'll want to know: aspirar}
\textbf{aspirar} — \textquotedblleft{}to vacuum, to hoover; to breathe in\textquotedblright{}.

\begin{quote}
Ao sábado aspiro a casa toda.
\end{quote}

\begin{quote}
\emph{On Saturdays I vacuum the whole house.}
\end{quote}

Also to aspire: \emph{aspirar a uma vida melhor}.

\begin{cousinweb}
\emph{aspirar} is a learned borrowing from Latin \emph{aspīrāre} \textquotedblleft{}to breathe on, to aspire\textquotedblright{} (EtymDB), built on \emph{spīrāre} \textquotedblleft{}to breathe\textquotedblright{}.
\end{cousinweb}

\begin{grammarlens}[title={What you've built}]
One more word for telling, explaining and describing at B1.
\end{grammarlens}

\subsection*{Your turn}
\begin{itemize}
\raggedright
  \item \emph{Say it:} \emph{aspirar}
  \item \emph{Say it:} \emph{aspirar}, once more
  \item \emph{Say it:} a sentence of your own with \emph{aspirar}
  \item \emph{Recall:} say the Portuguese for documentary, then the Portuguese for cast, then say \emph{aspirar} again
\end{itemize}

\begin{tcolorbox}[breakable,colback=teal!4,colframe=teal!35!black,title={Before you move on}]
What does aspirar mean? (\textquotedblleft{}To vacuum, to hoover; to breathe in\textquotedblright{}.) Say it once more.
\end{tcolorbox}
\section[o aspirador]{o aspirador — vacuum cleaner}

\label{lesson:PT-C528-aspirador}



\begin{quote}
Before the new one: say the Portuguese for cast, then the Portuguese for to vacuum.
\end{quote}

\subsection*{You'll want to know: o aspirador}
\textbf{o aspirador} — \textquotedblleft{}vacuum cleaner\textquotedblright{}.

\begin{quote}
O aspirador faz muito barulho.
\end{quote}

\begin{quote}
\emph{The vacuum cleaner makes a lot of noise.}
\end{quote}

Built on \emph{aspirar} with \emph{-dor}.

\begin{cousinweb}
\emph{aspirador} is built on \emph{aspirar} \textquotedblleft{}to breathe in, to suck in\textquotedblright{}, from Latin \emph{aspīrāre} (EtymDB).
\end{cousinweb}

\begin{grammarlens}[title={What you've built}]
One more word for telling, explaining and describing at B1.
\end{grammarlens}

\subsection*{Your turn}
\begin{itemize}
\raggedright
  \item \emph{Say it:} \emph{o aspirador}
  \item \emph{Say it:} \emph{o aspirador}, once more
  \item \emph{Say it:} a sentence of your own with \emph{o aspirador}
  \item \emph{Recall:} say the Portuguese for cast, then the Portuguese for to vacuum, then say \emph{o aspirador} again
\end{itemize}

\begin{tcolorbox}[breakable,colback=teal!4,colframe=teal!35!black,title={Before you move on}]
What does o aspirador mean? (\textquotedblleft{}Vacuum cleaner\textquotedblright{}.) Say it once more.
\end{tcolorbox}
\section[varrer]{varrer — to sweep}

\label{lesson:PT-C528-varrer}



\begin{quote}
Before the new one: say the Portuguese for to vacuum, then the Portuguese for vacuum cleaner.
\end{quote}

\subsection*{You'll want to know: varrer}
\textbf{varrer} — \textquotedblleft{}to sweep\textquotedblright{}.

\begin{quote}
Varri a cozinha depois do jantar.
\end{quote}

\begin{quote}
\emph{I swept the kitchen after dinner.}
\end{quote}

\emph{Varrer o chão} — to sweep the floor.

\begin{cousinweb}
\emph{varrer} is inherited, through Old Portuguese, from Latin \emph{verrere} \textquotedblleft{}to sweep\textquotedblright{} (EtymDB).
\end{cousinweb}

\begin{grammarlens}[title={What you've built}]
One more word for telling, explaining and describing at B1.
\end{grammarlens}

\subsection*{Your turn}
\begin{itemize}
\raggedright
  \item \emph{Say it:} \emph{varrer}
  \item \emph{Say it:} \emph{varrer}, once more
  \item \emph{Say it:} a sentence of your own with \emph{varrer}
  \item \emph{Recall:} say the Portuguese for to vacuum, then the Portuguese for vacuum cleaner, then say \emph{varrer} again
\end{itemize}

\begin{tcolorbox}[breakable,colback=teal!4,colframe=teal!35!black,title={Before you move on}]
What does varrer mean? (\textquotedblleft{}To sweep\textquotedblright{}.) Say it once more.
\end{tcolorbox}
\section[esfregar]{esfregar — to scrub, to rub}

\label{lesson:PT-C528-esfregar}



\begin{quote}
Before the new one: say the Portuguese for vacuum cleaner, then the Portuguese for to sweep.
\end{quote}

\subsection*{You'll want to know: esfregar}
\textbf{esfregar} — \textquotedblleft{}to scrub, to rub\textquotedblright{}.

\begin{quote}
Esfreguei o chão da casa de banho.
\end{quote}

\begin{quote}
\emph{I scrubbed the bathroom floor.}
\end{quote}

\emph{Esfregar os olhos} — to rub your eyes.

\begin{cousinweb}
\emph{esfregar} is inherited from Vulgar Latin \emph{exfricāre} (EtymDB), built on Latin \emph{fricāre} \textquotedblleft{}to rub\textquotedblright{}.
\end{cousinweb}

\begin{grammarlens}[title={What you've built}]
One more word for telling, explaining and describing at B1.
\end{grammarlens}

\subsection*{Your turn}
\begin{itemize}
\raggedright
  \item \emph{Say it:} \emph{esfregar}
  \item \emph{Say it:} \emph{esfregar}, once more
  \item \emph{Say it:} a sentence of your own with \emph{esfregar}
  \item \emph{Recall:} say the Portuguese for vacuum cleaner, then the Portuguese for to sweep, then say \emph{esfregar} again
\end{itemize}

\begin{tcolorbox}[breakable,colback=teal!4,colframe=teal!35!black,title={Before you move on}]
What does esfregar mean? (\textquotedblleft{}To scrub, to rub\textquotedblright{}.) Say it once more.
\end{tcolorbox}
\section[a esfregona]{a esfregona — mop}

\label{lesson:PT-C528-esfregona}



\begin{quote}
Before the new one: say the Portuguese for to sweep, then the Portuguese for to scrub.
\end{quote}

\subsection*{You'll want to know: a esfregona}
\textbf{a esfregona} — \textquotedblleft{}mop\textquotedblright{}.

\begin{quote}
Passei a esfregona no chão da cozinha.
\end{quote}

\begin{quote}
\emph{I mopped the kitchen floor.}
\end{quote}

Built on \emph{esfregar}. \emph{Passar a esfregona} — to mop.

\begin{cousinweb}
\emph{esfregona} is built on \emph{esfregar} \textquotedblleft{}to scrub\textquotedblright{}, inherited from Vulgar Latin \emph{exfricāre}, from Latin \emph{fricāre} \textquotedblleft{}to rub\textquotedblright{} (EtymDB).
\end{cousinweb}

\begin{grammarlens}[title={What you've built}]
One more word for telling, explaining and describing at B1.
\end{grammarlens}

\subsection*{Your turn}
\begin{itemize}
\raggedright
  \item \emph{Say it:} \emph{a esfregona}
  \item \emph{Say it:} \emph{a esfregona}, once more
  \item \emph{Say it:} a sentence of your own with \emph{a esfregona}
  \item \emph{Recall:} say the Portuguese for to sweep, then the Portuguese for to scrub, then say \emph{a esfregona} again
\end{itemize}

\begin{tcolorbox}[breakable,colback=teal!4,colframe=teal!35!black,title={Before you move on}]
What does a esfregona mean? (\textquotedblleft{}Mop\textquotedblright{}.) Say it once more.
\end{tcolorbox}
